%! TEX root = main.tex
\subsection{RegLez25}\linkdest{rl25}

\begin{frame}[allowframebreaks]{Reg Lez 25}
\begin{enumerate}
\item Lun 10/02/2025 lezione: Illustrazione del programma del corso e della modalit\'a d'esame. Caratteristiche generali delle stelle compatte e meccanismo di genesi. Propriet\'a generali delle pulsar: pulsar shape, periodo (P) e sua variazione nel tempo (Pdot). Interpretazione delle pulsar come stelle di neutroni ruotanti.
\item Ven 14/02/2025 lezione: Modello di dipolo magnetico ruotante delle pulsar: stima dell'intensit\'a del campo magnetico delle pulsar dalla misura di P e Pdot; braking index.
\item Lun 17/02/2025 lezione: Modello di dipolo magnetico ruotante delle pulsar (cont.): et\'a delle pulsar; et\'a della Crab pulsar; evoluzione delle pulsar nel piano P-Pdot; pulsar death line. Campi magnetici delle pulsar e cenni sulla loro origine. Effetti del decadimento nel tempo del campo magnetico sull'evoluzione delle pulsar, sulla determinazione dell'et\'a e sul braking index (caso di un campo magnetico che decade in modo esponenziale).
\item Ven 21/02/2025 lezione: Radiazione gravitazionale di una stella di neutroni ruotante. Confronto fra i casi di frenamento con braking index $n=3$ (dipolo magnetico) e $n=5$ (quadrupolo gravitazionale). Pulsar spin down limit: stima dell'eccentricit\'a della stella di neutroni associata alla Crab Nebula.
\item Lun 24/02/2025 lezione: Dipendenza dal tempo del momento di inerzia di una stella di neutroni: conseguenze sulla misura del campo magnetico, e dell'et\'a delle pulsars; braking index generalizzato $n(t)$. Equazione di stato di un gas di Fermi ideale non-relativistico.
\item Lun 03/03/2025 lezione: Equazione di stato di un gas di Fermi ideale relativistico; limite non-relativistico e ultra-relativistico. Equazione di stato per un ''core'' stellare degenere. Correzioni elettrostatiche all'equazione di stato di un gas ideale: approssimazione di Wigner-Seitz.
\item Ven 07/03/2025 lezione: Equazioni di struttura stellare in gravitazione newtoniana. Equazione di stato politropica. Equazione di Lane-Emden; soluzioni dell'equazione di Lane-Emden.
\item Lun 10/03/2025 lezione: Stelle nane bianche. Massa limite di Chandrasekhar. Effetto sul valore della massa di Chandrasekhar delle correzioni elettrostatiche all'equazione di stato di un gas ideale di Fermi. Tensore metrico e tensore energia-impulso per un fluido perfetto in Relativit\'a Speciale. Equazioni per la struttura stellare in relativit\'a generale: metrica in uno spazio-tempo statico e isotropo.
\item Lun 17/03/2025 lezione: Equazioni per la struttura stellare in relativit\'a generale (cont.): red-shift gravitazionale; tensore energia-impulso per un fluido perfetto. Equazione di Tolman-Oppenheimer-Volkov per l'equilibrio idrostatico: ruolo della EOS nella risoluzione numerica delle equazioni di TOV; Limite di Oppenheimer-Volkoff per la massa di una stella di neutroni.
\item Ven 21/03/2025 lezione: Equazioni per la struttura stellare in relativit\'a generale (cont.): massa gravitazionale, massa barionica e massa propria di una stella di neutroni; energia di legame di una stella di neutroni. Ruolo della relativit\'a speciale e della relativit\'a generale sulla massa limite di una stella compatta.
\item Lun 24/03/2025 lezione: Misura della massa di una stella di neutroni: sistemi binari, mass function, avanzamento del periastro, parametri post kepleriani. ''Upper bound'' (di Rhoades-Ruffini e Hartle-Sabbatini) sul valore della massa massima delle stelle di neutroni. Sistema di neutroni, protoni ed elettroni: potenziale chimico; condizione per il beta-equilibrio.
\item Ven 28/03/2025 lezione: Liquid drop model of the nucleus and semi-empirical formula of Bethe-Weizsaker for nuclear masses, Study of the stability conditions of isobaric nuclei by means of the semi-empirical mass formula. Nuclear Matter: saturation curve; properties at the empirical saturation point. Determination of Neutron star mass and number of neutrons from semi empirical formula, Equation of State (EoS). (Vishal Parmar)
\item Lun 31/03/2025 lezione: Brief Introduction of Mean-field theory. Meson picture of Nuclear forces, Relativistic Mean field theory (RMF): Sigma-omega model. Mean-field approximation, Solving the self-consistent equations of motion and plotting the eos for symmetric nuclear matter and Pure neutron matter. Beta-equilibrium and charge neutrality condition of neutron star matter. Scheme to solve the RMF equations for neutron star matter. Constraints on the EoS. (Vishal Parmar)
\item Ven 04/04/2025 lezione: Decadimento beta-inverso in un gas ideale di protoni ed elettroni ad alte densit\'a: soglia per i neutroni. Materia nucleare beta-stabile (caso di gas ideali): calcolo della frazione di protoni all'equilibrio.
\item Lun 07/04/2025 lezione: Materia nucleare beta-stabile (caso di gas ideali): calcolo della densit\'a di soglia per i muoni e per i pioni.
\item Ven 11/04/2025 lezione: Classificazione delle particelle: Adroni (barioni e mesoni), leptoni; cenni al modello a quark degli adroni. Stranezza ed adroni strani, produzione e decadimento; iperoni. Diagrammi degli adroni nel piano isospin-stranezza, ottetto e decupletto barionico. Calcolo della densit\'a di soglia per l'iperone Lambda in materia nucleare beta-stabile. Materia iperonica beta-stabile.
\item Lun 14/04/2025 lezione: Effetto dell'interazione forte sulla composizione della materia nucleare e iperonica beta-stabile; hyperon puzzle nelle stelle di neutroni (cenni). Nucleosintesi degli elementi oltre il ferro: r-process e s-process (cenni). La crosta delle stelle di neutroni: outer crust, neutron-drip density; inner crust. Transizione crosta-core: nuclear pasta phases (cenni). Materia nucleare asimmetrica, energia di simmetria. Materia nucleare beta-stabile; frazione di protoni all'equilibrio (caso con interazione nucleare); ruolo dell'energia di simmetria sulla massa limite e sul raggio della stella.
\item Lun 28/04/2025 lezione: Neutrino traping in materia beta-stabile. Protostelle di neutroni e loro evoluzione. Conseguanze del trapping dei neutrini sul concetto di massa limite di una stella di neutroni. Deconfinamento dei quark in materia densa. Strange Quark Matter (SQM) e relativa equazione di stato fenomenologica. Stelle di neutroni ibride. Ipotesi di Bodmer-Witten sulla assoluta stabilità della SQM. Strange stars: energia di legame e relazione massa-raggio delle strange stars.
\end{enumerate}
\end{frame}
